\documentclass[10pt,a4paper]{article}
\usepackage[T1]{fontenc}
\usepackage{lmodern,amsmath,amssymb,microtype,xcolor,geometry,hyperref}
\geometry{margin=19mm}
\definecolor{softgray}{gray}{0.94}
\hypersetup{hidelinks,pdftitle={SGA 2 Expose VIII Lemma 1.2 proof opening and Definition 1.3},pdfauthor={Interlanguage English/Germanic lane}}
\begin{document}
\begin{center}
  {\Large\bfseries SGA 2: Local Cohomology and Lefschetz Theorems}\par\smallskip
  {\large Expos\'e VIII: The Finiteness Theorem}\par
  \small Bounded source-aligned working unit --- 19 July 2026
\end{center}
\noindent\colorbox{softgray}{\parbox{0.95\linewidth}{\footnotesize\textbf{Authority.}
Corrected French lines 2561--2571; original printed pp.~85--86; physical
source-PDF p.~76; recomposed running p.~68. The original printed-page marker
changes to p.~86 within line 2563. This bounded unit has been independently
source-reviewed against the corrected French TeX and direct scan; cumulative
integration remains pending. Line 2572 is blank; Proposition 1.4 begins at
line 2573 and is excluded.}}

\emph{We shall consider only complexes whose differential has degree $+1$.}
By the hypothesis on $M$, there exists a projective resolution of $M$ of
finite length,
\[
  u\colon L^\bullet\longrightarrow M,
\]
where, moreover, the $L^p$ are finitely generated modules and $L^p=0$ if
$p\notin[-n,0]$. On the other hand, let
$v\colon M\longrightarrow I^\bullet$ be an injective resolution of $M$. I
claim that
\begin{equation}
  v\circ u\colon L^\bullet\longrightarrow I^\bullet
  \tag{1.1}
\end{equation}
is an injective resolution of $L^\bullet$. We must specify what this means.

\medskip
\noindent\textbf{Definition 1.3.} --- Let $X^\bullet$ be a complex of
$A$-modules. An \emph{injective resolution} of $X^\bullet$ is a homomorphism
of complexes
\[
  x\colon X^\bullet\longrightarrow CX^\bullet
\]
such that $CX^p$ is injective for every $p\in\mathbb Z$, and $x$ induces an
isomorphism in homology.
\end{document}
